% 原章第54—59页：附录5.A，式5.A.1—5.A.24。
\section*{附录 5.A\quad HSIC 的有偏与无偏估计量}
\phantomsection\addcontentsline{toc}{section}{附录 5.A：HSIC 的有偏与无偏估计量}
\markright{附录 5.A}

为使讨论完整，我们采用 Gretton（2015）的方法估计 HSIC。由式 (5.46) 得
\begin{equation}
\begin{aligned}
\mathrm{HSIC}(X,Y)
={}&\langle C_{XY}-\mu_x\otimes\mu_y,C_{XY}-\mu_x\otimes\mu_y\rangle\\
={}&\langle C_{XY},C_{XY}\rangle-2\langle C_{XY},\mu_x\otimes\mu_y\rangle
+\langle\mu_x\otimes\mu_y,\mu_x\otimes\mu_y\rangle.
\end{aligned}
\tag{5.A.1}
\end{equation}
根据未中心化协方差算子的定义，得到
\begin{equation}
\begin{aligned}
&\langle C_{XY},C_{XY}\rangle\\
&=\left\langle E_{p(x,y)}[k_x(\cdot,X)\otimes k_y(\cdot,Y)],
E_{p(x',y')}[k_x(\cdot,X')\otimes k_y(\cdot,Y')]\right\rangle\\
&=E_{p(x,y)}E_{p(x',y')}
\left[\langle k_x(\cdot,X)\otimes k_y(,Y),k_x(\cdot,X')\otimes k_y(\cdot,Y')\rangle\right]\\
&=E_{p(x,y)}E_{p(x',y')}
\left[\langle k_x(\cdot,X),k_x(,X')\rangle
\langle k_y(\cdot,Y),k_y(\cdot,Y')\rangle\right]\\
&=E_{p(x,y)}E_{p(x',y')}[k_x(X,X')k_y(Y,Y')].
\end{aligned}
\tag{5.A.2}
\end{equation}
利用期望的抽样公式，式 (5.A.2) 可化为
\begin{equation}
\begin{aligned}
\langle C_{XY},C_{XY}\rangle
&=\frac1{n^2}\sum_{i=1}^n\sum_{j=1}^nk_x(x_i,x_j)k_y(y_i,y_j)\\
&=\frac1{n^2}\sum_{i=1}^n\sum_{j=1}^nk_x(x_i,x_j)k_y(y_j,y_i)\\
&=\frac1{n^2}\tr(K_xK_y),
\end{aligned}
\tag{5.A.3}
\end{equation}
其中，$K_x=(k_x(x_i,x_j))_{n\times n}$，$K_y=(k_y(y_i,y_j))_{n\times n}$。

类似地，有
\begin{equation}
\begin{aligned}
&\langle C_{XY},\mu_x\otimes\mu_y\rangle\\
&=\left\langle E_{p(x,y)}\left[k_x(\cdot,X)\otimes k_y(\cdot,Y),
\left(E_{p(x)}[k_x(\cdot,X')]\right)\right.\right.\\
&\hspace{12em}\left.\left.\otimes\left(E_{p(y)}[k_y(\cdot,Y')]\right)\right]\right\rangle\\
&=E_{p(x,y)}\left[E_{p(x)}[k_x(X,X')]E_{p(y)}[k_y(Y,Y')]\right].
\end{aligned}
\tag{5.A.4}
\end{equation}
其抽样估计量为
\begin{equation}
\begin{aligned}
\langle C_{XY},\mu_x\otimes\mu_y\rangle
&=\frac1{n^2}E_{p(x,y)}\left[\sum_{j=1}^nk_x(X,x_j)\sum_{l=1}^nk_y(Y,y_l)\right]\\
&=\frac1{n^3}\sum_{i=1}^n\sum_{j=1}^n\sum_{l=1}^nk_x(x_i,x_j)k_y(y_i,y_l)\\
&=\frac1{n^3}\sum_{j=1}^n\sum_{i=1}^n\sum_{l=1}^nk_x(x_i,x_j)k_y(y_i,y_l)\\
&=\frac1{n^3}1_n^TK_xK_y1_n\\
&=\frac1{n^3}\tr(1_n1_n^TK_xK_y).
\end{aligned}
\tag{5.A.5}
\end{equation}
最后，可以证明
\begin{equation}
\begin{aligned}
&\langle\mu_x\otimes\mu_y,\mu_x\otimes\mu_y\rangle\\
&=\langle\mu_x,\mu_x\rangle\langle\mu_y,\mu_y\rangle\\
&=\langle E_x[k_x(\cdot,X)],E_{x'}[k_x(\cdot,X')]\rangle
\langle E_y[k_y(\cdot,Y)],E_{y'}[k_y(\cdot,Y')]\rangle\\
&=E_x[E_{x'}[k_x(X,X')]]E_y[E_{y'}[k_y(Y,Y')]].
\end{aligned}
\tag{5.A.6}
\end{equation}
同样，其抽样估计量可计算为
\begin{equation}
\begin{aligned}
&\langle\mu_x\otimes\mu_y,\mu_x\otimes\mu_y\rangle\\
&=\frac1{n^2}\sum_{i=1}^n\sum_{j=1}^nk_x(x_i,x_j)
\frac1{n^2}\sum_{q=1}^n\sum_{l=1}^nk_y(y_q,y_l)\\
&=\frac1{n^4}1_n^TK_x1_n1_n^TK_y1_n\\
&=\frac1{n^4}\tr(1_n1_n^TK_x1_n1_n^TK_y).
\end{aligned}
\tag{5.A.7}
\end{equation}
代入式 (5.A.3)、(5.A.5) 和 (5.A.7)，得到
\begin{equation}
\begin{aligned}
\mathrm{HSIC}(X,Y)
={}&\frac1{n^2}\Bigl[\tr(K_xK_y)-\frac2n\tr(1_n1_n^TK_xK_y)\\
&\hspace{7em}+\frac1{n^2}\tr(1_n1_n^TK_x1_n1_n^TK_y)\Bigr]\\
={}&\frac1{n^2}(K_xHK_yH),
\end{aligned}
\tag{5.A.8}
\end{equation}
其中，$H=I_n-\frac1n1_n1_n^T$。

现在推导 HSIC 的无偏估计量。此前的计算没有关注 $(x,y)$ 与 $(x',y')$ 的实际抽样方式，而是完全将二者的抽样视为相互独立。然而，$(x,y)$ 与 $(x',y')$ 实际上应当采用不放回抽样（Gretton，2015）。因此，$\lVert C_{XY}\rVert_{\mathrm{HS}}^2$ 的无偏抽样估计量为
\begin{equation}
\begin{aligned}
\langle C_{XY},C_{XY}\rangle
&=E_{p(x,y)}E_{p(x',y'}[k_x(X,X')k_y(Y,Y')]\\
&=\frac1{n(n-1)}\sum_{i=1}^n\sum_{j\ne i}^nk_x(x_i,x_j)k_y(y_i,y_j).
\end{aligned}
\tag{5.A.9}
\end{equation}
为简化记号，用 $\mathrm i_q^n$ 表示从 $\{1,\ldots,n\}$ 中抽取的全部 $q$ 元组构成的集合，并令 $(n)_q=n(n-1)\cdots(n-p+1)$。

采用这些记号，式 (5.A.9) 可化为
\begin{equation}
\langle C_{XY},C_{XY}\rangle
=\frac1{(n)_2}\sum_{(i,j)\in\mathrm i_2^n}k_x(x_i,x_j)k_y(y_i,y_j).
\tag{5.A.10}
\end{equation}
类似地，得到
\begin{equation}
\begin{aligned}
&\langle C_{XY},\mu_x\otimes\mu_y\rangle\\
&=E_{p(x,y)}[E_{p(x)}[k_x(X,X')]E_{p(y)}[k_y(Y,Y')]]\\
&=\frac1{n(n-1)}\sum_{j=1}^n\sum_{l\ne j}^nE_{p(x,y)}k_x(X,x_j)k_y(Y,y_l)\\
&=\frac1{n(n-1)(n-2)}\sum_{i\ne j,l}^n\sum_{l\ne j}^n\sum_{j=1}^nk_x(x_i,x_j)k_y(y_i,y_l)\\
&=\frac1{(n)_3}\sum_{(i,j,l)\in\mathrm i_3^n}k_x(x_i,x_j)k_y(y_i,y_l),
\end{aligned}
\tag{5.A.11}
\end{equation}
以及
\begin{equation}
\begin{aligned}
&\langle\mu_x\otimes\mu_y,\mu_x\otimes\mu_y\rangle\\
&=E_x[E_{x'}[k_x(X,X')]]E_y[E_{y'}[k_y(Y,Y')]]\\
&=\frac1{n(n-1)}\sum_{i=1}^n\sum_{j\ne i}^nE_x[k_x(X,x_i)]E_y[k_y(Y,y_j)]\\
&=\frac1{n(n-1)(n-2)(n-3)}
\sum_{i=1}^n\sum_{i\ne j}^n\sum_{q\ne i,j}^n\sum_{l\ne i,j,q}^nk_x(x_q,x_i)k_y(y_l,y_j)\\
&=\frac1{(n)_4}\sum_{(i,j,q,l)\in\mathrm i_4^n}k_x(x_i,x_q)k_y(y_j,y_l).
\end{aligned}
\tag{5.A.12}
\end{equation}

尽管式 (5.A.10)—(5.A.12) 是无偏估计量，它们的计算时间却远大于有偏估计量。接下来，降低有偏估计量的计算时间。令 $\widehat K_x$ 与 $\widehat K_y$ 分别表示将 $K_x$ 与 $K_y$ 的对角元素替换为零后得到的矩阵。于是，式 (5.A.9) 可改写为
\begin{equation}
\begin{aligned}
\langle C_{XY},C_{XY}\rangle
&=\frac1{(n)_2}\sum_{i=1}^n\sum_{j\ne i}^nk_x(x_i,x_j)k_y(y_i,y_j)\\
&=\frac1{(n)_2}\sum_{i=1}^n\sum_{j=1}^n\widehat k_x(x_i,x_j)\widehat k_y(y_j,y_i)\\
&=\frac1{(n)_2}\tr(\widehat K_x\widehat K_y).
\end{aligned}
\tag{5.A.13}
\end{equation}
类似地，利用式 (5.A.11)，得到
\begin{equation}
\begin{aligned}
&\langle C_{XY},\mu_x\otimes\mu_y\rangle\\
&=\frac1{n(n-1)(n-2)}\sum_{i=1}^n\sum_{j\ne i}^n\sum_{l\ne j,i}^nk_x(x_i,x_j)k_y(y_i,y_l)\\
&=\frac1{(n)_3}\sum_{j=1}^n\sum_{l\ne j}^n\sum_{i\ne j,l}^nk_x(x_i,x_j)k_y(y_i,y_l)\\
&=\frac1{(n)_3}\sum_{j=1}^n\sum_{i\ne j}^n\sum_{l\ne j,i}^nk_x(x_j,x_i)k_y(y_i,y_l)\\
&=\frac1{(n)_3}\sum_{j=1}^n\sum_{l\ne j}^n\sum_{i\ne j,l}^nk_x(x_j,x_i)k_y(y_i,y_l)\\
&=\frac1{(n)_3}\Bigl[\sum_{j=1}^n\sum_{l=1}^n\sum_{i\ne j,l}^nk_x(x_j,x_i)k_y(y_i,y_l)\\
&\hspace{9em}-\sum_{j=1}^n\sum_{i\ne j}k_x(x_j,x_i)k_y(y_i,y_j)\Bigr].
\end{aligned}
\tag{5.A.14}
\end{equation}
可以证明（习题 4）
\begin{equation}
\sum_{j=1}^n\sum_{l=1}^n\sum_{i\ne(j,l)}^nk_x(x_j,x_i)k_y(y_i,y_l)
=1_n^T\widehat K_x\widehat K_y1_n,
\tag{5.A.15}
\end{equation}
以及
\begin{equation}
\sum_{j=1}^n\sum_{i\ne j}k_x(x_j,x_i)k_y(y_i,y_j)
=\tr(\widehat K_x\widehat K_y).
\tag{5.A.16}
\end{equation}
代入式 (5.A.15) 和 (5.A.16)，得到
\begin{equation}
\langle C_{XY},\mu_x\otimes\mu_y\rangle
=\frac1{(n)_3}\left[1_n^T\widehat K_x\widehat K_y1_n-\tr(\widehat K_x\widehat K_y)\right].
\tag{5.A.17}
\end{equation}
再次利用式 (5.A.12)，得到
\begin{equation}
\langle\mu_x\otimes\mu_y,\mu_x\otimes\mu_y\rangle
=\frac1{(n)_4}\sum_{(i,j,q,l)\in\mathrm i_4^n}k_x(x_i,x_q)k_y(y_j,y_l).
\tag{5.A.18}
\end{equation}
注意，下式左侧乘积中的各项可展开如下（Gretton，2015；习题 5）：
\begin{equation}
\begin{aligned}
&\sum_{(i,q)\in\mathrm i_2^n}k_x(x_i,x_q)\sum_{(j,l)\in\mathrm i_2^n}k_y(y_j,y_l)\\
&=\sum_{(i,q,j)\in\mathrm i_3^n}k_x(x_i,x_q)k_y(y_j,y_i)\\
&\quad+\sum_{(i,q,l)\in\mathrm i_3^n}k_x(x_i,x_q)k_y(y_q,y_l)\\
&\quad+\sum_{(i,j,l)\in\mathrm i_3^n}k_x(x_i,x_j)k_y(y_j,y_l)\\
&\quad+\sum_{(q,j,l)\in\mathrm i_3^n}k_x(x_l,x_q)k_y(y_j,y_l)\\
&\quad+\sum_{(i,j,q,l)\in\mathrm i_4^n}k_x(x_i,x_q)k_y(y_j,y_l)\\
&\quad+\sum_{(i,q)\in\mathrm i_2^n}k_x(x_i,x_q)k_y(y_i,y_q)\\
&\quad+\sum_{(i,q)\in\mathrm i_2^n}k_x(x_i,x_q)k_y(y_i,y_q).
\end{aligned}
\tag{5.A.19}
\end{equation}
显然，
\begin{equation}
\begin{aligned}
\sum_{(i,q)\in\mathrm i_2^n}k_x(x_i,x_q)&=1_n^T\widehat K_x1_n,\\
\sum_{(j,l)\in\mathrm i_2^n}k_y(y_j,y_l)&=1_n^T\widehat K_y1_n.
\end{aligned}
\tag{5.A.20}
\end{equation}
经过简单的代数运算，得到
\begin{equation}
\begin{aligned}
1_n^T\widehat K_x\widehat K_y1_n
&=\sum_{i=1}^n\sum_{j\ne i}^n\sum_{l\ne j}^nk_x(x_i,x_j)k_y(y_j,y_l)\\
&=\sum_{i=1}^n\sum_{j\ne i}^n\sum_{l\ne(i,j)}^nk_x(x_i,x_j)k_y(y_j,y_l)\\
&\quad+\sum_{i=1}^n\sum_{j\ne i}^nk_x(x_i,x_j)k_y(y_j,y_i)\\
&=\sum_{(i,j,l)\in\mathrm i_3^n}k_x(x_i,x_j)k_y(y_j,y_l)+\tr(\widehat K_x\widehat K_y).
\end{aligned}
\tag{5.A.21}
\end{equation}
采用类似论证，得到
\begin{equation}
\begin{aligned}
1_n^T\widehat K_x\widehat K_y1_n
&=\sum_{(i,q,j)\in\mathrm i_3^n}k_x(x_i,x_q)k_y(y_j,y_i)+\tr(\widehat K_x\widehat K_y),\\
1_n^T\widehat K_x\widehat K_y1_n
&=\sum_{(i,q,l)\in\mathrm i_3^n}k_x(x_i,x_q)k_y(y_q,y_l)+\tr(\widehat K_x\widehat K_y),\\
1_n^T\widehat K_x\widehat K_y1_n
&=\sum_{(q,j,l)\in\mathrm i_3^n}k_x(x_l,x_q)k_y(y_j,y_l)+\tr(\widehat K_x\widehat K_y).
\end{aligned}
\tag{5.A.22}
\end{equation}
将式 (5.A.20)—(5.A.22) 代入式 (5.A.19)，得到
\begin{equation}
\begin{aligned}
&\langle\mu_x\otimes\mu_y,\mu_x\otimes\mu_y\rangle\\
&=\frac1{(n)_4}\left[(1_n^T\widehat K_x1_n)(1_n^T\widehat K_y1_n)
-4\,1_n^T\widehat K_x\widehat K_y1_n+2\tr(\widehat K_x\widehat K_y)\right].
\end{aligned}
\tag{5.A.23}
\end{equation}
最后，将式 (5.A.13)、(5.A.17) 和 (5.A.23) 代入式 (5.A.1)，得到 HSIC 的无偏估计量：
\begin{equation}
\begin{aligned}
\mathrm{HSIC}_u(X,Y)
={}&\frac1{(n)_2}\tr(\widehat K_x\widehat K_y)
-\frac2{(n)_3}\left[1_n^T\widehat K_x\widehat K_y1_n-\tr(\widehat K_x\widehat K_y)\right]\\
&+\frac1{(n)_4}\Bigl[(1_n^T\widehat K_y1_n)(1_n^T\widehat K_y1_n)
-4\,1_n^T\widehat K_x\widehat K_y1_n\\
&\hspace{20em}+2\tr(\widehat K_x\widehat K_y)\Bigr]\\
={}&\left(\frac1{(n)_2}+\frac2{(n)_3}+\frac1{(n)_4}\right)\tr(\widehat K_x\widehat K_y)
-\left(\frac2{(n)_3}+\frac1{(n)_4}\right)1_n^T\widehat K_x\widehat K_y1\\
&+\frac1{(n)_4}(1_n^T\widehat K_x1_n)(1_n^T\widehat K_y1_n)\\
={}&\frac1{n(n-3)}\tr(\widehat K_x\widehat K_y)
-\frac2{n(n-2)(n-3)}1_n^T\widehat K_x\widehat K_y1\\
&+\frac1{n(n-1)(n-2)(n-3)}(1_n^T\widehat K_x1_n)(1_n^T\widehat K_y1_n)\\
={}&\frac1{n(n-3)}\Bigl[\tr(\widehat K_x\widehat K_y)
-\frac2{n-2}1_n^T\widehat K_x\widehat K_y1\\
&\hspace{4em}+\frac1{(n-1)(n-2)}(1_n^T\widehat K_x1_n)(1_n^T\widehat K_y1_n)\Bigr].
\end{aligned}
\tag{5.A.24}
\end{equation}
